\subsection{Process Ancestry Analysis}
The core innovation of the PRM Framework is its process ancestry analysis. When a system call is intercepted, the framework:

\begin{enumerate}
\item Extracts the command names of the current process, its parent, and grandparent
\item Compares this ancestry against defined security rules
\item Makes a security decision based on the matching rule
\end{enumerate}

This approach provides rich context for security decisions, allowing the framework to distinguish between identical operations initiated through different execution paths.

\subsection{PRM Table Rules}
PRM Table Rules define the security policies that the Verifier uses to determine whether a system call is legitimate. Each rule specifies:

\begin{enumerate}
\item Process pattern: What processes the rule applies to
\item Parent pattern: What parent processes are expected
\item Grandparent pattern: What grandparent processes are expected
\item Action: What action to take when a match is found (ACCEPT, REJECT, REDIRECT, RETURN, DEBUG)
\item Hook type: What system call types the rule applies to
\end{enumerate}

\subsection{Performance Optimizations}
The PRM Framework incorporates several performance optimizations:

\begin{enumerate}
\item \textbf{Efficient Caching}: An LRU cache stores previous security decisions to avoid redundant processing
\item \textbf{Fast Path Processing}: Early decisions for common operations
\item \textbf{Minimal Processing}: Early returns and optimized memory usage
\end{enumerate}

\section{Evaluation}
\label{sec:evaluation}

\subsection{Evaluation Methodology}
To evaluate the PRM Framework, we conducted a series of experiments comparing its performance and security characteristics with existing solutions. Our methodology consisted of:

\begin{enumerate}
\item \textbf{Benchmark Selection}: We selected standard HPC benchmarks (LINPACK, STREAM, IOR, IMB-MPI) that represent typical computational, memory, I/O, and communication workloads in HPC environments.

\item \textbf{Test Environment}: All benchmarks were performed on a 16-node cluster with Intel Xeon Gold 6248 processors (20 cores/40 threads per node), 128GB RAM per node, and InfiniBand HDR interconnect. Each node ran CentOS 8.3 with Linux kernel 5.4.

\item \textbf{Security Configurations}: We tested each security solution with configurations that provided comparable security guarantees:
   \begin{itemize}
   \item SELinux: Targeted policy in enforcing mode
   \item AppArmor: Default profiles with container confinement
   \item seccomp: Docker default profile plus custom rules
   \item PRM Framework: Process ancestry rules matching the protection level of other solutions
   \end{itemize}

\item \textbf{Measurement Protocol}: Each benchmark was run 10 times, and we report the average values with standard deviation. System call overhead was measured using the \texttt{strace} tool with timing enabled.
\end{enumerate}

\subsection{Performance Comparison}
Table \ref{tab:performance} provides a quantitative comparison of performance overhead between PRM Framework and other security solutions, based on our measurements and published literature.

\begin{table}[htbp]
\caption{Performance overhead comparison between different security solutions}
\label{tab:performance}
\centering
\begin{tabular}{|l|c|c|c|c|c|}
\hline
\textbf{Security Solution} & \textbf{System Call} & \textbf{File Access} & \textbf{Network} & \textbf{Memory} & \textbf{Container} \\
 & \textbf{Overhead} & \textbf{Overhead} & \textbf{Operation} & \textbf{Operation} & \textbf{Startup} \\
 &  &  & \textbf{Overhead} & \textbf{Overhead} & \textbf{Overhead} \\
\hline
\textbf{PRM Framework} & 0.5-2\% ($\pm$0.3\%) & 1-3\% ($\pm$0.5\%) & 0.8-2.5\% ($\pm$0.4\%) & 0.3-1.5\% ($\pm$0.2\%) & 2-5\% ($\pm$0.8\%) \\
\hline
\textbf{SELinux} & 5-10\% ($\pm$1.2\%)[12] & 7-15\% ($\pm$2.1\%)[12] & 4-9\% ($\pm$1.5\%)[12] & 3-7\% ($\pm$0.9\%)[12] & 10-20\% ($\pm$3.2\%)[13] \\
\hline
\textbf{AppArmor} & 3-7\% ($\pm$0.8\%)[5] & 5-10\% ($\pm$1.4\%)[5] & 3-7\% ($\pm$1.1\%)[5] & 2-5\% ($\pm$0.7\%)[5] & 5-15\% ($\pm$2.5\%)[13] \\
\hline
\textbf{seccomp} & 0.3-1\% ($\pm$0.2\%)[14] & N/A & 0.5-2\% ($\pm$0.3\%)[14] & 0.3-1\% ($\pm$0.2\%)[14] & 2-5\% ($\pm$0.6\%)[13] \\
\hline
\end{tabular}
\begin{tablenotes}
\small
\item Values show mean overhead ranges with standard deviation in parentheses. Lower percentages indicate better performance (less overhead). Sources: Reshetova et al. \cite{reshetova2016security}, Goldberg et al. \cite{goldberg2019performance}, Bui \cite{bui2015analysis}, Grattafiori \cite{grattafiori2016understanding}, and our measurements for PRM Framework.
\end{tablenotes}
\end{table}

Table \ref{tab:features} presents feature comparison between PRM Framework and other security solutions based on our analysis and published literature.

\begin{table}[htbp]
\caption{Feature comparison between different security solutions}
\label{tab:features}
\centering
\begin{tabular}{|l|c|c|c|c|}
\hline
\textbf{Feature} & \textbf{PRM Framework} & \textbf{SELinux} & \textbf{AppArmor} & \textbf{seccomp} \\
\hline
\textbf{Implementation} & eBPF programs & LSM hooks & LSM hooks & seccomp-bpf \\
\hline
\textbf{Context Awareness} & Process ancestry + & Label-based & Path-based & System call \\
 & UID transitions &  &  & numbers + args \\
\hline
\textbf{Dynamic Updates} & Yes (with & Requires & Requires & Limited \\
 & recompilation) & policy reload & profile reload &  \\
\hline
\textbf{Container Support} & Excellent & Good & Good & Good \\
\hline
\textbf{Security Granularity} & Process ancestry & Resource labels & Path-based & System call filtering \\
\hline
\textbf{Learning Curve} & Moderate & Steep[15] & Moderate[15] & Low[15] \\
\hline
\end{tabular}
\begin{tablenotes}
\small
\item Source: Schreuders et al. \cite{schreuders2011empowering} and our analysis.
\end{tablenotes}
\end{table}